\documentclass[12pt]{article}
\def\activityid{F07-X-v2}\usepackage[letterpaper,margin=.65in,top=.60in,bottom=.65in]{geometry}
\usepackage[T1]{fontenc}\usepackage[default]{sourcesanspro}
\usepackage{microtype,tikz,array,tabularx,fancyhdr,amsmath}\usepackage[hidelinks]{hyperref}
\usetikzlibrary{arrows.meta,calc}
\setlength{\parindent}{0pt}\setlength{\parskip}{7pt}\setlength{\headheight}{0pt}\setlength{\footskip}{26pt}\setlength{\emergencystretch}{2em}
\pagestyle{fancy}\fancyhf{}\renewcommand{\headrulewidth}{0pt}\renewcommand{\footrulewidth}{.4pt}
\fancyfoot[L]{\scriptsize Bellingham Math Circle / Week 7 / \activityid}\fancyfoot[R]{\scriptsize \thepage}
\definecolor{ink}{gray}{.12}\definecolor{soft}{gray}{.95}\color{ink}
\newcommand{\PageTitle}[2]{{\fontsize{10}{12}\selectfont\bfseries WEEK 7\quad GAME LAB\hfill #1\par}\vspace{3pt}{\fontsize{26}{29}\selectfont\bfseries #2\par}\vspace{2pt}}
\newcommand{\NameLine}{{\small Name: \rule{2.65in}{.4pt}\hfill Date: \rule{1.35in}{.4pt}\par}\vspace{4pt}}
\newcommand{\Task}[2]{\par\vspace{3pt}{\large\bfseries #1}\par #2\par}
\newcommand{\Subhead}[1]{\par\vspace{3pt}{\large\bfseries #1}\par}
\newcommand{\RuleBox}[1]{\par\begingroup\setlength{\fboxsep}{8pt}\colorbox{soft}{\parbox{\dimexpr\linewidth-16pt\relax}{#1}}\endgroup\par}
\newcommand{\WriteBox}[2]{\par\textbf{#1}\par\vspace{2pt}\begin{tikzpicture}\draw[black!40,line width=.5pt] (0,0) rectangle (\dimexpr\linewidth-.5pt\relax,#2);\end{tikzpicture}\par}
\newcommand{\StudentFooter}{\vfill{\small Build first. Explain by pointing, talking, or drawing.}\par}
\newcommand{\Code}[1]{\texttt{#1}}

\hypersetup{pdftitle={Week 7: extra-grades-6-7}}
\begin{document}
\PageTitle{EXTRA / GRADES 6--7}{Two winning piles can make a loss}\NameLine
\RuleBox{Use moves 1, 3, 4. Now there are \textbf{two piles}; a turn changes exactly one pile. The player who takes the very last counter wins.}
\Task{1. W and L forget too much.}{Play from piles (1,1), then (1,4). Each individual pile is W. Are both combined games W? Explain your best play.}
\Task{2. Give each pile a richer label.}{Let \(g(0)=0\). For a larger pile, list the \(g\)-values reachable in one move; give the pile the \textbf{smallest nonnegative integer missing} from that list. For example, \(g(1)=1\) and \(g(2)=0\).}\begin{center}\renewcommand{\arraystretch}{1.9}\begin{tabular}{c|ccccccccccc}Pile&0&1&2&3&4&5&6&7&8&9&10\\\hline$g$&0&1&0&&&&&&&&\end{tabular}\end{center}
\Task{3. Test an equality rule.}{Conjecture: two piles are losing together exactly when their \(g\)-values are equal. Check (1,3), (4,6), and (5,6). Find a winning move wherever possible.}
\WriteBox{Checks and moves:}{.7in}
\Task{4. Prove the rule for two piles.}{Why can a move never keep equal labels equal? If labels differ, why can you move from the larger label to the smaller one? Use the ``smallest missing'' definition.}
\WriteBox{The two parts of my proof:}{.8in}
\textbf{Beyond this page:} Three piles require binary XOR, not ordinary addition or pairwise equality. These labels are Sprague--Grundy values.\StudentFooter
\end{document}
